\chapter{Problémafelvetés és motiváció}

\section{Miért nem elég egy hagyományos statisztikai oldal?}

A legtöbb, játékosoknak szánt statisztikai felület egyetlen mérkőzésre, egyetlen rangsorra vagy néhány alapmutatóra szűkíti le az értelmezést. Ez önmagában hasznos lehet, de nem válaszolja meg azt a kérdést, amely a projekt egészét végigkísérte: vajon a mérkőzések mögött kirajzolódik-e egy tartós társas szerkezet, és ha igen, ez a szerkezet hogyan kapcsolódik a teljesítményhez?

A \emph{Premade Graph} két egymástól nem független problémával szembesül. Az egyik adatmérnöki jellegű: a nyers Riot API válaszokból valahogy olyan adatkorpuszt kell felépíteni, ahol a játékosok azonosítása, a meccsek tárolása és a gráfprojekció újrafuttatható módon működik. A másik inkább kutatási: mit lehet egyáltalán állítani ezekből az adatokból, ami nem egyszerűen látványos, hanem valóban alátámasztható is.

\section{A projekt természete}

A kiinduló adathalmaz nem egyszerű játékoslista, hanem ismétlődő interakciós háló. A rendszer ugyanazt a két játékost nemcsak egyszeri ellenfelekként látja, hanem potenciális visszatérő szövetségesként, visszatérő ellenfélként, klasztertagként, hídcsúcsként vagy egy nagyobb globális hálózat részeként is.

Ez több szempontból vizsgálható:

\begin{itemize}
    \item a League of Legends mérkőzésadatai elég gazdagok ahhoz, hogy valódi teljesítményprofilok legyenek számíthatók;
    \item a visszatérő együttjátszás és egymás elleni előfordulás gráfelméleti szempontból értelmezhető kapcsolatokat hoz létre;
    \item a szövetségesi és ellenfél-jelleg együttese miatt a projekt nem csak súlyozott, hanem többféle kapcsolatprojekcióként is vizsgálható;
    \item a rendszer implementációs oldalon is érdekes, mert több nyelv és futtatókörnyezet közötti felelősségmegosztást kell megindokolni.
\end{itemize}

\section{A központi szakdolgozati állítás}

A dolgozat fő tézise az, hogy \emph{a valós mérkőzéstörténetből felépített League of Legends játékosháló egyszerre kezelhető mint asszociatív játékosgráf, mint teljesítményprofilokkal dúsított adatstruktúra, és mint rendszertechnikai szempontból komoly, többnyelvű analitikai platform}.

Ez három fő irányban bontható ki, egy kísérleti kiterjesztéssel kiegészülve. A hálózati irány azt vizsgálja, hogy az ismétlődő együttjátszás és egymás elleni előfordulás mérhető szerkezeti mintázatot alkot-e: a \textit{flexset} adatkészleten ez leginkább asszociatív core-periphery szerkezetként mutatkozik meg. A metrikai irány a \textit{feedscore} és \textit{opscore} pontszámokra épül. A rendszertechnikai irány azt indokolja, miért szükséges Python, SQLite, Node/Express, Rust és React együttese. A kísérleti irány a Tribal NeuroSim szimulációra vonatkozik.

\section{Hatókör és korlátok}

A projekt viszonylag nagy rendszer, és könnyen csúszhat abba a hibába, hogy minden meglévő ötlet kész eredményként szerepeljen. A dolgozat igyekszik ezt elkerülni. Ami ténylegesen fut és ellenőrizhető, az kész állapotként szerepel. Amit a fejlesztési tervek jövőbeli célként fogalmaznak meg, az nem szerepel kész feature-ként.

\section{A dolgozat szerkezete}

A dolgozat fejlődési ívet követ: az adatok eredetétől a rendszer egyre összetettebb rétegein át jut el a két adatkészlet empirikus összehasonlításáig, majd a szimulációs kísérletig.

Az első fejezetek a kiindulópontot rögzítik: a problémafelvetést, a szakirodalmi hátteret és a Riot API-ra épülő adatforrásokat. Ezt követi a rendszerarchitektúra bemutatása, majd az \textit{opscore} teljesítménymetrika, a gráfépítés és klaszterezés mechanizmusa, a Rust futtatókörnyezet és az analitikai modulok. A két adatkészlet (\textit{flexset} és \textit{soloq}) empirikus összehasonlítása a rendszer legkonkrétabb kutatási eredménye. A dolgozat utolsó egységei az adatbázis-engineering alprojektet és a Tribal NeuroSim szimulációt tárgyalják.
